% GENERATED FILE. Edit canonical lessons, then run npm run generate:books.
% canonical-source-hash: fnv1a64:e7a310f8aaa9ecec
% canonical-lessons: ZH-C127-zhexie, ZH-W127-xie-few, ZH-C127-naxie, ZH-C127-naxie-which, ZH-C127-yixie, ZH-C127-zheyang

\chapter{These and Those}
\label{ch:zh-these-and-those}
\hlchaptermodality{127}

\begin{chapteropening}
\textbf{By the end of this chapter:} \emph{I can say these, those, which ones, some, and like this.}

Complete the last lesson of chapter 127: I can say these, those, which ones, some, and like this.
\end{chapteropening}

\section[\texorpdfstring{\zh{这些} (zhèxiē)}{这些 (zhèxiē)}]{\zh{这些} (zhèxiē) — these}

\label{lesson:ZH-C127-zhexie}



\begin{quote}
Before the new one: say the Mandarin for a moment ago, then the Mandarin for already.
\end{quote}

\subsection*{You'll want to know: \zh{这些}}
\textbf{\zh{这些}} — \emph{zhèxiē} — \textquotedblleft{}these\textquotedblright{}.

\begin{quote}
\textbf{\zh{这些是我以前的书}} \emph{zhèxiē shì wǒ yǐqián de shū} — These are my old books.
\end{quote}

The new character, in pieces: \zh{些} — on top a four-stroke piece and a two-stroke piece side by side, and under them the two bars of \zh{二}. The next lesson puts it in your hand.

\begin{cousinweb}
\textbf{\zh{这些}} (\emph{zhèxiē}) is built from \zh{这} 'this' and \zh{些} 'a few'. The second is inherited from Middle Chinese \emph{sjae} (Guangyun, in Baxter's transcription), glossed 'a little'.
\end{cousinweb}

\begin{grammarlens}[title={What you've built}]
One more word for pointing things out.
\end{grammarlens}

\subsection*{Your turn}
\begin{itemize}
\raggedright
  \item \emph{Say it:} \emph{zhèxiē}
  \item \emph{Say it:} \emph{zhèxiē}, once more
  \item \emph{Say it:} \emph{zhèxiē shì wǒ yǐqián de shū} — These are my old books.
  \item \emph{Recall:} say the Mandarin for a moment ago, then the Mandarin for already, then say \emph{zhèxiē} again
\end{itemize}

\begin{tcolorbox}[breakable,colback=teal!4,colframe=teal!35!black,title={Before you move on}]
What does \zh{这些} mean? (\textquotedblleft{}These\textquotedblright{}.) Say it once more.
\end{tcolorbox}
\section[\texorpdfstring{\zh{些} (xiē)}{些 (xiē)}]{\zh{些} — two small pieces over \zh{二}}

\label{lesson:ZH-W127-xie-few}



\begin{quote}
The word first, then the character inside it.

\begin{itemize}
\raggedright
  \item \emph{Recall:} say \emph{zhèxiē} — \textbf{R1}, one lesson back, before it decays
\end{itemize}

\begin{quote}
\zh{这} \zh{二}
\end{quote}

\emph{Read it:} each one aloud, with nothing said first — \emph{zhè}, \emph{èr}
\end{quote}

\subsection*{Script — two small pieces over \zh{二}}
\hlblockfigure{\detokenize{figures/ZH-W127-xie-few-filmstrip.pdf}}{How \zh{些} is written, stroke by stroke}

\begin{quote}
\zh{些}
\end{quote}

\textbf{Three pieces: two side by side on top, then \zh{二} underneath.}

\textbf{The top left piece, four strokes.} A short vertical, a short horizontal from its middle to the right, a shorter vertical to the left of the first, then a stroke along the foot that rises to the right.

\textbf{The top right piece, two strokes.} A short left-falling stroke, then a vertical that bends to the right along the bottom and hooks up.

\textbf{\zh{二}, two strokes}: the short bar, then the long one beneath it.

4 + 2 + 2 = eight.

\subsection*{Your turn}
\begin{itemize}
\raggedright
  \item \emph{Write it:} \zh{些}, all eight, with the model in front of you
  \item \emph{Cover:} write it again from memory, and count the pieces, not the strokes
\end{itemize}

\begin{tcolorbox}[breakable,colback=teal!4,colframe=teal!35!black,title={Before you move on}]
How many pieces in \zh{些}? (\textbf{Three}.) Which character you know sits at the bottom? (\textbf{\zh{二}}.) How does the sixth stroke end? (\textbf{It bends right along the bottom and hooks up}.)

Source: \href{https://raw.githubusercontent.com/chanind/hanzi-writer-data/68d10a4b21150cae5e1ebbd223eed289cf32d90c/data/\%E4\%BA\%9B.json}{Hanzi Writer Data \zh{些}, pinned stroke paths and medians}.
\end{tcolorbox}
\section[\texorpdfstring{\zh{那些} (nàxiē)}{那些 (nàxiē)}]{\zh{那些} (nàxiē) — those}

\label{lesson:ZH-C127-naxie}



\begin{quote}
Before the new one: say the Mandarin for already, then the Mandarin for these.
\end{quote}

\subsection*{You'll want to know: \zh{那些}}
\textbf{\zh{那些}} — \emph{nàxiē} — \textquotedblleft{}those\textquotedblright{}.

\begin{quote}
\textbf{\zh{那些人是我以前的同学}} \emph{nàxiē rén shì wǒ yǐqián de tóngxué} — Those people are my old classmates.
\end{quote}

\begin{cousinweb}
\textbf{\zh{那些}} (\emph{nàxiē}) is built from \zh{那} 'that' and \zh{些} 'a few', both inherited from Middle Chinese, \emph{na} and \emph{sjae} (Guangyun, in Baxter's transcription). The dictionary glosses the first 'how, what', not 'that', and the second 'a little'.
\end{cousinweb}

\begin{grammarlens}[title={What you've built}]
One more word for pointing things out.
\end{grammarlens}

\subsection*{Your turn}
\begin{itemize}
\raggedright
  \item \emph{Say it:} \emph{nàxiē}
  \item \emph{Say it:} \emph{nàxiē}, once more
  \item \emph{Say it:} \emph{nàxiē rén shì wǒ yǐqián de tóngxué} — Those people are my old classmates.
  \item \emph{Recall:} say the Mandarin for already, then the Mandarin for these, then say \emph{nàxiē} again
\end{itemize}

\begin{tcolorbox}[breakable,colback=teal!4,colframe=teal!35!black,title={Before you move on}]
What does \zh{那些} mean? (\textquotedblleft{}Those\textquotedblright{}.) Say it once more.
\end{tcolorbox}
\section[\texorpdfstring{\zh{哪些} (nǎxiē)}{哪些 (nǎxiē)}]{\zh{哪些} (nǎxiē) — which ones}

\label{lesson:ZH-C127-naxie-which}



\begin{quote}
Before the new one: say the Mandarin for these, then the Mandarin for those.
\end{quote}

\subsection*{You'll want to know: \zh{哪些}}
\textbf{\zh{哪些}} — \emph{nǎxiē} — \textquotedblleft{}which ones\textquotedblright{}.

\begin{quote}
\textbf{\zh{你喜欢哪些书}} \emph{nǐ xǐhuan nǎxiē shū} — Which books do you like?
\end{quote}

\begin{cousinweb}
\textbf{\zh{哪些}} (\emph{nǎxiē}) is built from \zh{哪} 'which' and \zh{些} 'a few'. The second is inherited from Middle Chinese \emph{sjae} (Guangyun, in Baxter's transcription); no dictionary we use records an older form of the first, so its history is not settled.
\end{cousinweb}

\begin{grammarlens}[title={What you've built}]
One more word for pointing things out.
\end{grammarlens}

\subsection*{Your turn}
\begin{itemize}
\raggedright
  \item \emph{Say it:} \emph{nǎxiē}
  \item \emph{Say it:} \emph{nǎxiē}, once more
  \item \emph{Say it:} \emph{nǐ xǐhuan nǎxiē shū} — Which books do you like?
  \item \emph{Recall:} say the Mandarin for these, then the Mandarin for those, then say \emph{nǎxiē} again
\end{itemize}

\begin{tcolorbox}[breakable,colback=teal!4,colframe=teal!35!black,title={Before you move on}]
What does \zh{哪些} mean? (\textquotedblleft{}Which ones\textquotedblright{}.) Say it once more.
\end{tcolorbox}
\section[\texorpdfstring{\zh{一些} (yìxiē)}{一些 (yìxiē)}]{\zh{一些} (yìxiē) — some; a few}

\label{lesson:ZH-C127-yixie}



\begin{quote}
Before the new one: say the Mandarin for those, then the Mandarin for which ones.
\end{quote}

\subsection*{You'll want to know: \zh{一些}}
\textbf{\zh{一些}} — \emph{yìxiē} — \textquotedblleft{}some; a few\textquotedblright{}.

\begin{quote}
\textbf{\zh{我想喝一些水}} \emph{wǒ xiǎng hē yìxiē shuǐ} — I'd like to drink some water.
\end{quote}

The tone of \zh{一} changes here: yī is said yì before the first tone of xiē.

\begin{cousinweb}
\textbf{\zh{一些}} (\emph{yìxiē}) is built from \zh{一} 'one' and \zh{些} 'a few', both inherited from Middle Chinese, \emph{'jit} and \emph{sjae} (Guangyun, in Baxter's transcription).
\end{cousinweb}

\begin{grammarlens}[title={What you've built}]
One more word for pointing things out.
\end{grammarlens}

\subsection*{Your turn}
\begin{itemize}
\raggedright
  \item \emph{Say it:} \emph{yìxiē}
  \item \emph{Say it:} \emph{yìxiē}, once more
  \item \emph{Say it:} \emph{wǒ xiǎng hē yìxiē shuǐ} — I'd like to drink some water.
  \item \emph{Recall:} say the Mandarin for those, then the Mandarin for which ones, then say \emph{yìxiē} again
\end{itemize}

\begin{tcolorbox}[breakable,colback=teal!4,colframe=teal!35!black,title={Before you move on}]
What does \zh{一些} mean? (\textquotedblleft{}Some; a few\textquotedblright{}.) Say it once more.
\end{tcolorbox}
\section[zhèyàng]{zhèyàng — like this; this way}

\label{lesson:ZH-C127-zheyang}



\begin{quote}
Before the new one: say the Mandarin for which ones, then the Mandarin for a few.
\end{quote}

\subsection*{You'll want to know: zhèyàng}
\emph{zhèyàng} — \textquotedblleft{}like this; this way\textquotedblright{}.

\begin{quote}
\emph{nǐ kěyǐ zhèyàng zuò} — You can do it this way.
\end{quote}

Its characters come later; for now it is a word for the ear.

\begin{cousinweb}
\emph{zhèyàng} is built from \zh{这} 'this' and \emph{yàng} 'a pattern', the \emph{yàng} of \emph{yíyàng}. No dictionary we use records an older form of either part in these senses, so the history of the pair is not settled.
\end{cousinweb}

\begin{grammarlens}[title={What you've built}]
One more word for pointing things out.
\end{grammarlens}

\subsection*{Your turn}
\begin{itemize}
\raggedright
  \item \emph{Say it:} \emph{zhèyàng}
  \item \emph{Say it:} \emph{zhèyàng}, once more
  \item \emph{Say it:} \emph{nǐ kěyǐ zhèyàng zuò} — You can do it this way.
  \item \emph{Recall:} say the Mandarin for which ones, then the Mandarin for a few, then say \emph{zhèyàng} again
\end{itemize}

\begin{tcolorbox}[breakable,colback=teal!4,colframe=teal!35!black,title={Before you move on}]
What does zhèyàng mean? (\textquotedblleft{}Like this; this way\textquotedblright{}.) Say it once more.
\end{tcolorbox}
